% Carl Friedrich Gauss, Werke IX, printed 304–306, Article [5] through case I.
\par\bigskip
\begin{center}
[5.]\\[4pt]
Condition equations of the second kind.
\end{center}

[Once the triangle misclosures have been eliminated by adjustment, the
condition equations arising from the ratios of the sides are set up.

For the quadrilateral $7.\ 8.\ 10.\,.\,.\,.\,.\,.\ 9$, with 9 chosen as the central point:
\begin{center}
\small
\setlength{\tabcolsep}{3pt}
\renewcommand{\arraystretch}{1.25}
\begin{tabular}{@{}r r@{: }r r r l@{}}
\shortstack{in the\\triangle}&\multicolumn{2}{c}{adjusted angles}&$\frac13$ excess&$\log\sin$&correction\quad]\\
7 at&$7$&$66^\circ\ 1'\ 19''{,}151$&$-2''{,}269$&$9{,}960\,8021.983$&$+9.36\{(7.9)-(7.8)\}$\\
&$8$&$66\ \ 39\ \ 58{,}666$&$-2{,}269$&$9{,}962\,9416.342$&$+9.08\{(8.7)-(8.9)\}$\\[4pt]
$9$&$10$&$22\ \ 59\ \ 17{,}428$&$-0{,}794$&$9{,}591\,6628.409$&$+49.63\{(10.7)-(10.9)\}$\\
&$7$&$10\ \ 27\ \ 3{,}217$&$-0{,}794$&$9{,}258\,6108.427$&$+114.15\{(7.9)-(7.10)\}$\\
\end{tabular}
\end{center}
\newpage
\noindent\textbf{ON NETWORK ADJUSTMENT.}\hfill305
\begin{center}
\small
\setlength{\tabcolsep}{3pt}
\renewcommand{\arraystretch}{1.25}
\begin{tabular}{@{}r r@{: }r r r l@{}}
$\bigl[$\shortstack{in the\\triangle}&\multicolumn{2}{c}{adjusted angles}&$\frac13$ excess&$\log\sin$&correction\quad]\\
10 at&$8$&$23^\circ\ 11'\ 52''{,}167$&$-1''{,}415$&$9{,}595\,3867.516$&$+49.13\{(8.9)-(8.10)\}$\\
&$10$&$57\ \ 33\ \ 19{,}333$&$-1{,}415$&$9{,}926\,2943.910$&$+13.39\{(10.8)-(10.9)\}$\\
&&&&$9{,}997\,8605.641$&\\
&&&&$0{,}333\,0519.982$&\\
&&&&$9{,}669\,0923.606$&\\
\cline{5-5}
&&&&$+49.229$&\\
&&&&\textnormal{(previously $+57$) [*)]}&\\
\end{tabular}
\end{center}

[Hence the side equation for $7.\ 8.\ 10.\,.\,.\,.\,.\,.\ 9$ is:]
\[
\begin{aligned}
0={}&+49{,}229-9{,}36(7.8)-104{,}79(7.9)+114{,}15(7.10)\\
&-9{,}08(8.7)+58{,}21(8.9)-49{,}13(8.10)\\
&+49{,}63(10.7)-13{,}39(10.8)-36{,}24(10.9).
\end{aligned}
\]

[This equation is transformed as follows. If the adjustment
of a figure requires satisfying a number of angle equations with zero constant
terms and just one side equation, this can be done by adjusting
\textls[120]{a single} equation. This equation is obtained by multiplying
the angle equations by certain factors
and adding them to the side equation. These factors are to be chosen
so that the sum of the squares of the coefficients in the transformed
side equation is a minimum. Thus, take the already adjusted angle equations
for triangles 7, 9, 10:
\[
\begin{aligned}
+(7.9)-(9.7)+(9.8)-(8.9)+(8.7)-(7.8)&=0\\
+(7.9)-(9.7)+(9.10)-(10.9)+(10.7)-(7.10)&=0\\
+(8.9)-(9.8)+(9.10)-(10.9)+(10.8)-(8.10)&=0,
\end{aligned}
\]
Multiplying them by $x,y,z$ and adding them to the side equation above
first gives:
\[
\begin{aligned}
0={}&+49{,}229-(9{,}36+x)(7.8)-(104{,}79-x-y)(7.9)+(114{,}15-y)(7.10)\\
&-(9{,}08-x)(8.7)+(58{,}21-x+z)(8.9)-(49{,}13+z)(8.10)\\
&-(x+y)(9.7)+(x-z)(9.8)+(y+z)(9.10)\\
&+(49{,}63+y)(10.7)-(13{,}39-z)(10.8)-(36{,}24+y+z)(10.9),
\end{aligned}
\]
\par\smallskip
\noindent\rule{0.25\textwidth}{0.4pt}\par
{\footnotesize
[*) The value in parentheses, $+57$, is obtained by calculating the side equation using
the observed angles in Art. 3.]\par}
\vfill
\noindent\textsc{IX.}\hfill39
\newpage
\noindent306\hfill\textbf{POSTHUMOUS PAPERS.}\par\medskip
\noindent and if we now minimize the sum of the squares of the coefficients
of the corrections, then
\[
\begin{aligned}
6x+2y-2z&=+162{,}7\\
2x+6y+2z&=+133{,}1\\
-2x+2y+6z&=-130{,}2
\end{aligned}
\]
must hold.] We obtain
\[
\begin{aligned}
x&=+7{,}76\\
y&=+29{,}21\\
z&=-28{,}85,
\end{aligned}
\]
and hence the corrected condition equation:
\[
\begin{aligned}
\mathrm{I.}\quad 0={}&+49{,}229-17{,}12(7.8)-67{,}82(7.9)+84{,}94(7.10)\\
&-1{,}32(8.7)+21{,}60(8.9)-20{,}28(8.10)\\
&-36{,}97(9.7)+36{,}61(9.8)+0{,}36(9.10)\\
&+78{,}84(10.7)-42{,}24(10.8)-36{,}60(10.9).
\end{aligned}
\]

[In what follows, only the original and transformed side equations
are given. The heading always identifies the figure to which the side equation
belongs and the triangles whose angle equations are used
to transform it. The triangle numbers are written in place of the factors
$x,y,z$.]
